\section*{Bab \ref{ch:g3:continents-and-oceans} --- Bola Dunia: Benua dan Samudra}

\begin{solution}{exo:g3:continents-and-oceans:1}
Asia, Afrika, Amerika Utara, Amerika Selatan, Antarktika, Eropa, dan
Australia. Yang terbesar adalah Asia; yang terkecil adalah Australia.
\end{solution}

\begin{solution}{exo:g3:continents-and-oceans:2}
\omterm{def:g2:water-on-earth:ocean}{Samudra} Pasifik, Atlantik, Hindia, Selatan, dan Arktik. Yang terbesar
adalah \omterm{def:g2:water-on-earth:ocean}{Samudra} Pasifik; yang terkecil adalah \omterm{def:g2:water-on-earth:ocean}{Samudra} Arktik.
\end{solution}

\begin{solution}{exo:g3:continents-and-oceans:3}
\omterm{def:g3:continents-and-oceans:poles}{Kutub Utara} ada di bagian atas \omterm{def:g3:continents-and-oceans:globe}{bola dunia}; \omterm{def:g3:continents-and-oceans:poles}{Kutub Selatan} ada di bagian
bawah; \omterm{def:g3:continents-and-oceans:equator}{khatulistiwa} melingkari bagian tengahnya.
\end{solution}

\begin{solution}{exo:g3:continents-and-oceans:4}
Eropa: utara. Australia: selatan. Antarktika: selatan. Amerika Utara:
utara.
\end{solution}

\begin{solution}{exo:g3:continents-and-oceans:5}
\omterm{def:g2:water-on-earth:ocean}{Samudra} Arktik terletak di sekitar \omterm{def:g3:continents-and-oceans:poles}{Kutub Utara}. \omterm{def:g2:water-on-earth:ocean}{Samudra} Selatan terletak di
sekitar Antarktika.
\end{solution}

\begin{solution}{exo:g3:continents-and-oceans:6}
Dari Afrika ke Amerika Selatan: \omterm{def:g2:water-on-earth:ocean}{Samudra} Atlantik. Dari Australia ke arah
timur menuju Amerika Selatan: \omterm{def:g2:water-on-earth:ocean}{Samudra} Pasifik. Dari Afrika ke Australia:
\omterm{def:g2:water-on-earth:ocean}{Samudra} Hindia.
\end{solution}

\begin{solution}{exo:g3:continents-and-oceans:7}
Sam tinggal di dekat \omterm{def:g3:continents-and-oceans:equator}{khatulistiwa}, yang panas sepanjang tahun. Bo tinggal
di dekat \omterm{def:g3:continents-and-oceans:poles}{kutub}, yang dingin sepanjang tahun.
\end{solution}

\begin{solution}{exo:g3:continents-and-oceans:8}
$10 - 7 = 3$: jarinya mendarat di daratan 3 kali. Inilah yang bisa kita
duga: sekitar 71 persegi dari 100 persegi Bumi berupa air, jadi kira-kira
7 kali dari 10 kali jarinya berhenti di \omterm{def:g2:water-on-earth:ocean}{laut}. (Dengan hanya 10 putaran,
Lena bisa juga mendapat 6 atau 8: hasil hitungannya berubah dari satu
permainan ke permainan berikutnya.)
\end{solution}

\begin{solution}{exo:g3:continents-and-oceans:9}
Untuk menggambar \omterm{def:g3:continents-and-oceans:world-map}{peta dunia}, \omterm{def:g3:continents-and-oceans:globe}{bola dunia} dipotong sepanjang garis yang
melewati tengah \omterm{def:g2:water-on-earth:ocean}{Samudra} Pasifik lalu dibuat datar. Separuh \omterm{def:g2:water-on-earth:ocean}{Samudra} Pasifik
berada di tepi kiri \omterm{def:g2:drawing-to-map:map}{peta}, separuh lainnya di tepi kanan. Di Bumi yang
sebenarnya hanya ada satu \omterm{def:g2:water-on-earth:ocean}{Samudra} Pasifik: kedua tepi \omterm{def:g2:drawing-to-map:map}{peta} itu adalah
garis yang sama pada \omterm{def:g3:continents-and-oceans:globe}{bola dunia}.
\end{solution}

\begin{solution}{exo:g3:continents-and-oceans:10}
Tom salah. Dari yang terbesar sampai yang terkecil, \omterm{def:g3:continents-and-oceans:continent}{benua-benua} itu adalah
Asia, Afrika, Amerika Utara, Amerika Selatan, Antarktika, Eropa, dan
Australia: Antarktika hanya urutan kelima, jauh lebih kecil daripada Asia.
\omterm{def:g3:continents-and-oceans:world-map}{Peta dunia} mengecohnya karena \omterm{def:g2:drawing-to-map:map}{peta} datar sangat meregangkan daratan di
dekat \omterm{def:g3:continents-and-oceans:poles}{kutub}. Antarktika, yang terletak di sekitar \omterm{def:g3:continents-and-oceans:poles}{Kutub Selatan}, terentang
menjadi pita panjang di sepanjang seluruh bagian bawah \omterm{def:g2:drawing-to-map:map}{peta}. Pada \omterm{def:g3:continents-and-oceans:globe}{bola
dunia}, Antarktika adalah daratan yang bundar, dan jauh lebih kecil daripada
Asia.
\end{solution}
